% Feedback:
% 1. Years of Experience — maybe consider two pages. 
%    Write as much as you need to. Thanks you.
% 2. The first page is therefore prime real estate. 
% 3. Skills section shorter. 
% 4. Don't bother with portfolio. You have years of experience. 
% 5. 

%%%%%%%%%%%%%%%%%%%%%%%%%%%%%%%%%%%%%%%%%%%%%%%%%%%%%%%%%%%%%%%%%%%%%%%%
%%%%%%%%%%%%%%%%%%%%%% Simple LaTeX CV Template %%%%%%%%%%%%%%%%%%%%%%%%
%%%%%%%%%%%%%%%%%%%%%%%%%%%%%%%%%%%%%%%%%%%%%%%%%%%%%%%%%%%%%%%%%%%%%%%%

%%%%%%%%%%%%%%%%%%%%%%%%%%%%%%%%%%%%%%%%%%%%%%%%%%%%%%%%%%%%%%%%%%%%%%%%
%% NOTE: If you find that it says                                     %%
%%                                                                    %%
%%                           1 of ??                                  %%
%%                                                                    %%
%% at the bottom of your first page, this means that the AUX file     %%
%% was not available when you ran LaTeX on this source. Simply RERUN  %%
%% LaTeX to get the ``??'' replaced with the number of the last page  %%
%% of the document. The AUX file will be generated on the first run   %%
%% of LaTeX and used on the second run to fill in all of the          %%
%% references.                                                        %%
%%%%%%%%%%%%%%%%%%%%%%%%%%%%%%%%%%%%%%%%%%%%%%%%%%%%%%%%%%%%%%%%%%%%%%%%

%%%%%%%%%%%%%%%%%%%%%%%%%%%% Document Setup %%%%%%%%%%%%%%%%%%%%%%%%%%%%

% Don't like 10pt? Try 11pt or 12pt
\documentclass[10pt]{article}

% The automated optical recognition software used to digitize resume
% information works best with fonts that do not have serifs. This
% command uses a sans serif font throughout. Uncomment both lines (or at
% least the second) to restore a Roman font (i.e., a font with serifs).
%\usepackage{times}
%\renewcommand{\familydefault}{\sfdefault}

% This is a helpful package that puts math inside length specifications
\usepackage{calc}
\usepackage{comment}
\usepackage{graphicx}
% \usepackage{fontawesome}

% Simpler bibsection for CV sections
% (thanks to natbib for inspiration)
\makeatletter
\newlength{\bibhang}
\setlength{\bibhang}{1em} %1em}
\newlength{\bibsep}
 {\@listi \global\bibsep\itemsep \global\advance\bibsep by\parsep}
\newenvironment{bibsection}%
		{\begin{enumerate}{}{%
%        {\begin{list}{}{%
	   \setlength{\leftmargin}{\bibhang}%
	   \setlength{\itemindent}{-\leftmargin}%
	   \setlength{\itemsep}{\bibsep}%
	   \setlength{\parsep}{\z@}%
		\setlength{\partopsep}{0pt}%
		\setlength{\topsep}{0pt}}}
		{\end{enumerate}\vspace{-.6\baselineskip}}
%        {\end{list}\vspace{-.6\baselineskip}}
\makeatother

% Layout: Puts the section titles on left side of page
\reversemarginpar

%
%         PAPER SIZE, PAGE NUMBER, AND DOCUMENT LAYOUT NOTES:
%
% The next \usepackage line changes the layout for CV style section
% headings as marginal notes. It also sets up the paper size as either
% letter or A4. By default, letter was used. If A4 paper is desired,
% comment out the letterpaper lines and uncomment the a4paper lines.
%
% As you can see, the margin widths and section title widths can be
% easily adjusted.
%
% ALSO: Notice that the includefoot option can be commented OUT in order
% to put the PAGE NUMBER *IN* the bottom margin. This will make the
% effective text area larger.
%
% IF YOU WISH TO REMOVE THE ``of LASTPAGE'' next to each page number,
% see the note about the +LP and -LP lines below. Comment out the +LP
% and uncomment the -LP.
%
% IF YOU WISH TO REMOVE PAGE NUMBERS, be sure that the includefoot line
% is uncommented and ALSO uncomment the \pagestyle{empty} a few lines
% below.
%

%% Use these lines for letter-sized paper
%\usepackage[paper=letterpaper,
%            %includefoot, % Uncomment to put page number above margin
%            marginparwidth=1.2in,     % Length of section titles
%            marginparsep=.05in,       % Space between titles and text
%            margin=1in,               % 1 inch margins
%            includemp]{geometry}

% Use these lines for A4-sized paper
\usepackage[paper=a4paper,
			%includefoot, % Uncomment to put page number above margin
			marginparwidth=0mm,    % Length of section titles
			marginparsep=0mm,       % Space between titles and text
			margin=18mm,              % 25mm margins
			includemp]{geometry}

%% More layout: Get rid of indenting throughout entire document
\setlength{\parindent}{0in}

\usepackage[shortlabels]{enumitem}

%% Reference the last page in the page number
%
% NOTE: comment the +LP line and uncomment the -LP line to have page
%       numbers without the ``of ##'' last page reference)
%
% NOTE: uncomment the \pagestyle{empty} line to get rid of all page
%       numbers (make sure includefoot is commented out above)
%
\usepackage{fancyhdr,lastpage}
\pagestyle{fancy}
%\pagestyle{empty}      % Uncomment this to get rid of page numbers
\fancyhf{}\renewcommand{\headrulewidth}{0pt}
\fancyfootoffset{\marginparsep+\marginparwidth}
\newlength{\footpageshift}
\setlength{\footpageshift}
		  {0.5\textwidth+0.5\marginparsep+0.5\marginparwidth-2in}
\lfoot{\hspace{\footpageshift}%
	   \parbox{4in}{\, \hfill %
					\arabic{page} of \protect\pageref*{LastPage} % +LP
%                    \arabic{page}                               % -LP
					\hfill \,}}

% Finally, give us PDF bookmarks
\usepackage{color,hyperref}
\definecolor{darkblue}{rgb}{0.0,0.0,0.3}
\hypersetup{colorlinks,breaklinks,
			linkcolor=darkblue,urlcolor=darkblue,
			anchorcolor=darkblue,citecolor=darkblue}

%%%%%%%%%%%%%%%%%%%%%%%% End Document Setup %%%%%%%%%%%%%%%%%%%%%%%%%%%%


%%%%%%%%%%%%%%%%%%%%%%%%%%% Helper Commands %%%%%%%%%%%%%%%%%%%%%%%%%%%%

% The title (name) with a horizontal rule under it
% (optional argument typesets an object right-justified across from name
%  as well)
%
% Usage: \makeheading{name}
%        OR
%        \makeheading[right_object]{name}
%
% Place at top of document. It should be the first thing.
% If ``right_object'' is provided in the square-braced optional
% argument, it will be right justified on the same line as ``name'' at
% the top of the CV. For example:
%
%       \makeheading[\emph{Curriculum vitae}]{Your Name}
%
% will put an emphasized ``Curriculum vitae'' at the top of the document
% as a title. Likewise, a picture could be included:
%
%   \makeheading[\includegraphics[height=1.5in]{my_picutre}]{Your Name}
%
% the picture will be flush right across from the name.
\newcommand{\makeheading}[8][]{
	\hspace*{-\marginparsep minus \marginparwidth}%
	\begin{minipage}[t]{\textwidth+\marginparwidth+\marginparsep}%
		\begin{tabular}[t]{@{}l}
		{\hspace{-8pt} \Large \bfseries#2}\\
		\\
		{\normalsize \emph{#3}}
		\end{tabular}
		\hfill
		\begin{tabular}[t]{|l}
		#4  \\
		#5  \\
		#6 \\
		
		\end{tabular}
		\vspace{1pt}
		\\[-0.2\baselineskip]%
		
		\rule{\columnwidth}{1pt}%
	\end{minipage}
}

% The section headings
%
% Usage: \section{section name}
% \renewcommand{\section}[1]{\pagebreak[3]%
%     \hyphenpenalty=10000%
%     \vspace{1.8\baselineskip}%
%     \phantomsection\addcontentsline{toc}{section}{#1}%
%     \noindent\llap{\scshape\smash{\parbox[t]{\marginparwidth}{\raggedright #1}}}%
%     \vspace{-\baselineskip}\par}

\renewcommand{\section}[1]{
	\pagebreak[3]%
	\hyphenpenalty=10000%
	\vspace{\baselineskip}%
	\phantomsection\addcontentsline{toc}{section}{#1}%
	\noindent{\scshape #1}%
	\vspace{\baselineskip}\par
}

\newcommand{\sectioncont}[1]{
	\pagebreak[3]%
	\hyphenpenalty=10000%
	\vspace{\baselineskip}%
	\noindent{\scshape #1} (continued)% 
	\vspace{\baselineskip}\par
}

% An itemize-style list with lots of space between items
\newenvironment{outerlist}[1][\enskip\textbullet]{
\begin{itemize}[#1,leftmargin=*]}{\end{itemize}%
\vspace{-.6\baselineskip}
}

% An environment IDENTICAL to outerlist that has better pre-list spacing
% when used as the first thing in a \section
\newenvironment{lonelist}[1][\enskip\textbullet]{
	\begin{list}{#1}{%
	\setlength{\partopsep}{0pt}%
	\setlength{\topsep}{0pt}}
}{\end{list}\vspace{-.6\baselineskip}}

% An itemize-style list with little space between items
\newenvironment{innerlist}[1][\enskip\textbullet]{
	\begin{itemize}[#1,leftmargin=*,parsep=0pt,itemsep=1pt,topsep=2pt,partopsep=0pt]
}{\end{itemize}}

\newenvironment{innerlist2}[1][\enskip\textasteriskcentered]{
	\begin{itemize}[#1,leftmargin=*,parsep=0pt,itemsep=1pt,topsep=2pt,partopsep=0pt]
}{\end{itemize}}

% An environment IDENTICAL to innerlist that has better pre-list spacing
% when used as the first thing in a \section
\newenvironment{loneinnerlist}[1][\enskip\textbullet]%
		{\begin{itemize}[#1,leftmargin=*,parsep=0pt,itemsep=0pt,topsep=0pt,partopsep=0pt]}
		{\end{itemize}\vspace{-.6\baselineskip}}

% To add some paragraph space between lines.
% This also tells LaTeX to preferably break a page on one of these gaps
% if there is a needed pagebreak nearby.
\newcommand{\blankline}{\quad\pagebreak[3]}
\newcommand{\halfblankline}{\quad\vspace{-0.5\baselineskip}\pagebreak[3]}

% Uses hyperref to link DOI
\newcommand\doilink[1]{\href{http://dx.doi.org/#1}{#1}}
\newcommand\doi[1]{doi:\doilink{#1}}

% For \url{SOME_URL}, links SOME_URL to the url SOME_URL
\providecommand*\url[1]{\href{#1}{#1}}
% Same as above, but pretty-prints SOME_URL in teletype fixed-width font
\renewcommand*\url[1]{\href{#1}{\texttt{#1}}}

% For \email{ADDRESS}, links ADDRESS to the url mailto:ADDRESS
\providecommand*\email[1]{\href{mailto:#1}{#1}}
% Same as above, but pretty-prints ADDRESS in teletype fixed-width font
%\renewcommand*\email[1]{\href{mailto:#1}{\texttt{#1}}}

\newcommand{\nohref}[2]{#2}

%\providecommand\BibTeX{{\rm B\kern-.05em{\sc i\kern-.025em b}\kern-.08em
%    T\kern-.1667em\lower.7ex\hbox{E}\kern-.125emX}}
%\providecommand\BibTeX{{\rm B\kern-.05em{\sc i\kern-.025em b}\kern-.08em
%    \TeX}}
\providecommand\BibTeX{{B\kern-.05em{\sc i\kern-.025em b}\kern-.08em
	\TeX}}
\providecommand\Matlab{\textsc{Matlab}}

%%%%%%%%%%%%%%%%%%%%%%%% End Helper Commands %%%%%%%%%%%%%%%%%%%%%%%%%%%

%%%%%%%%%%%%%%%%%%%%%%%%% Begin CV Document %%%%%%%%%%%%%%%%%%%%%%%%%%%%
\newlength{\rcollength}\setlength{\rcollength}{1.2in}%


\begin{document}
\makeheading{
		\href{https://joshua.prettyman.me/}{Joshua Prettyman, Ph.D.}
}
{Data Scientist, Mathematics Ph.D., Python Developer}
% {Python Developer, Data Scientist, Mathematics Ph.D.}
{\email{joshua@prettyman.me}}
% {\href{https://github.com/DrPrettyman}{github.com/DrPrettyman}}
{\href{https://www.linkedin.com/in/joshuaprettyman/}{linkedin.com/in/joshuaprettyman}}
% {{\href{tel:+44(0)7470885347}{tel: 07470 885 347}}}
% {{\href{tel:+34711042515}{(+34) 711 042 151}}}
{\href{https://joshua.prettyman.me/}{joshua.prettyman.me}}



%\section{Objective}

%Insert text here if you want to
%\begin{innerlist}
%\item More information and auxiliary documents can be found at\\\url{http://www.tedpavlic.com/facjobsearch/}
%\end{innerlist}



\section{Personal Statement}

I developed Blink SEO's database and internal software from scratch
and oversaw its evolution into the ML-powered SaaS product for e-commerce
\href{https://www.macaroni.works}{\textit{Macaroni Software}}, 
boasting \href{https://blinkseo.co.uk/blogs/case-studies/how-macaroni-made-us-20x-more-productive}{a $20\times$ increase in SEO productivity}. 
I created the role after being hired as a one-man data/dev team with the brief:
`use data to improve processes'. 
The team has grown and I have been refining the product in an agile environment.
% I am looking for a role in which I can continue to solve problems and 
% develop software using python.
I am looking for a new role where I can implement my data science skills and 
continue my professional development. 

% \vspace{.05in}

% I have  experience lecturing upto Masters level, presenting at international conferences 
% and \href{https://www.linkedin.com/in/joshuaprettyman/details/publications/}{publishing in respected journals}, so I'm used 
% to communicating technical concepts to a variety of audiences. 

\vspace{-6pt}
\begin{center}
		\rule{0.5\columnwidth}{1pt}
\end{center}
\vspace{-10pt}

% \section{Tech Stack}
\section{Data Science Skills and Tech Stack}
\vspace{-6pt}
\begin{innerlist}
	% \item Deployment to remote UNIX server using Docker and GitHub.
	\item Regular use of Python libraries: ScikitLearn, Numpy, Pandas, NLTK, MatPlotLib, Plotly, etc..
	\item Python development: Data Science / ML, async, package management, REST APIs, etc..
	\item Full stack development with Python, Google Cloud Platform, BigQuery, PostgreSQL and JavaScript.
	\item Huggingface transformers and LLMs, ollama LLM, OpenAI Assistant API.
	% \item I have experience producing visualisations for public apps, business reports, 
	% 	dashboards and academic publications.
	\item \href{https://www.codecademy.com/learn/paths/data-science}{CodeCademy} pro 
		\textit{Data Scientist} course:
		ScikitLearn, TensorFlow, PyTorch, NLP and data visualisation.
	\item Experience C++, shell scripting, Matlab, Retool, 
		data engineering, dashboards, scientific computing. 
\end{innerlist}

\vspace{-6pt}
\begin{center}
		\rule{0.5\columnwidth}{1pt}
\end{center}
\vspace{-10pt}

\section{Professional Experience}
\setlength\tabcolsep{0pt}

\textbf{Data Scientist} — \href{https://www.blinkseo.co.uk}{Blink SEO} / \href{https://www.macaroni.works}{Macaroni Software} \hfill \emph{Nov. 2021 to Dec. 2024} \\
\vspace{-6pt}
\begin{innerlist}
\item I delivered Macaroni Software, allowing SEO teams to deliver 
	a year's work \href{https://www.linkedin.com/posts/sam-wright-17b6ab6_shopify-seo-activity-7170336529146441729-JGDn?utm_source=combined_share_message&utm_medium=member_desktop}{in a single month}.
\item I created this role by working with the team to replicate 
	and then improve their (spreadsheet-based) processes with continual feedback and deployment.
\item Macaroni does the (typically manual) data-engineering for SEO by aggregating data from various sources, 
	streamed daily (or on request) to a \emph{BigQuery} database via the python backend (\textit{GCP Compute Engine}: Linux/Debian) which utilises 
	various REST APIs and preprocessing steps.
\item Client onboarding, data imports, data analysis tasks and ML-assisted recommendations can be triggered 
	in the app and run asynchronously in the backend using a (\textit{PostgreSQL}-based) job-queue system which I wrote for this purpose.
\item Quantitative data are used for various visualisations in the app, powered by \textit{Plotly}.
\item Text data (on-page content and search terms) are processed and combined with 
	quantitative data (clicks, impressions, etc.) to surface actionable recommendations 
	for new content and site taxonomy. This uses NLP and Clustering algorithms (\textit{NLTK}, \textit{huggingface}, \textit{ScikitLearn}) but also 
	a deep understanding of SEO, necessitating regular exchanges with the SEO delivery team 
	and frequent iterations. 
\item I also integrated an LLM (\textit{huggingface}) to generate ready-to-go suggestions 
	for content gaps, allowing customers to see immediate impact. 
\item Besides conceiving of and working on Macaroni, I also did frequent, 
	one-off data tasks for the SEO delivery and management teams including 
	data engineering, content generation, and providing analyses and 
	visualisations for clients and investors. 
\end{innerlist}

\vspace{14pt}
\textbf{Data Science Researcher} — \href{http://www.npl.co.uk}{National Physical Laboratory (NPL)} \hfill \emph{Sep. 2015 to Feb. 2021} \\
\vspace{-6pt}
\begin{innerlist}
	\item I completed my doctoral research in industrial collaboration 
		with Dr. Valerie Livina at NPL.
	\item I presented data and results for stake-holder meetings at NLP, 
		as well as in published articles and at international conferences. 
	\item My research developed methods to detect tipping points in dynamical systems, 
		working with large time series datasets 
		and building stochastic models to test hypotheses. 
	\item Code written in \textit{MatLab} and \textit{Python}. Visualisations produced 
		using \textit{MatPlotLib} and \textit{XMGrace}.
	\item Signals in hourly sea-level pressure were used to produce early warning signals 
		for approaching tropical storms. This required the cleaning, interpolating and 
		correct interpretation of large meteorological datasets downloaded (raw) from 
		public sources.
	
	% \item Work involved analysing large datasets and coding in Matlab and Python.
\end{innerlist}


\pagebreak

\sectioncont{Professional Experience}

\textbf{Associate Lecturer} — \href{https://www.shu.ac.uk/about-us/academic-departments/engineering-and-mathematics}{Sheffield Hallam University} \hfill \emph{Sep. 2017 to Jul. 2019} \\
\vspace{-6pt}
\begin{innerlist}
	\item I taught mathematics, statistics and computing courses from Foundation Degree to Masters.
	\item Responsibilities included lesson preparation, course planning, marking, leading tutorials. 
\end{innerlist} 
\vspace{10pt}

\textbf{Informatics developer (internship)} — \href{https://www.metoffice.gov.uk/research/foundation/informatics-lab/what-we-do}{UK Met Office Informatics Lab} \hfill \emph{Jun. to Aug. 2017} \\
\vspace{-6pt}
\begin{innerlist}
	\item A three-month project assessing the viability of using live data from 
		amateur meteorologists to augment official observation data. 
	\item This involved accessing internal and external data via APIs using Java, 
		then performing data cleaning and analyses in Python.
	\item I worked as part of a team and participated in daily stand-ups.
\end{innerlist} 



%\vspace{.1in}
%\textbf{Informatics developer} — \href{https://www.metoffice.gov.uk/research/foundation/informatics-lab/what-we-do}{UK Met Office Informatics Lab} \hfill \emph{2017}
%\begin{innerlist}
%	\item 3 Month project involved consolidating and testing observation data. 
%%	\item Worked as part of a team: daily stand-ups, using Slack to coordinate.
%	\item Working in Python, accessing internal and external data via APIs.
%\end{innerlist}

%\textbf{Teaching Assistant} — University of Reading \hfill \emph{2014 - 2015}
%\vspace{.05in}
%	\begin{innerlist}
%		\item Tutorial support and teaching cover on undergraduate courses whilst undertaking my own Masters studies.
%		\item Responsibilities such as marking and exam invigilator.
%	\end{innerlist}



\vspace{1pt}
\begin{center}
		\rule{0.5\columnwidth}{1pt}
\end{center}
\vspace{-10pt}

% \pagebreak


\section{Academic background}
\setlength\tabcolsep{0pt}

\textbf{Ph.D. Mathematics} — \textbf{University of Reading}\\
\vspace{-2pt}
\hspace{2pt} \textbf{Thesis:} \href{http://centaur.reading.ac.uk/98364/1/23022044_Prettyman_Thesis_Joshua\%20Prettyman.pdf}{\emph{Tipping Points and Early Warning Signals with Applications to Geophysical Data}}.
\vspace{4pt}
\begin{innerlist}
	\item My thesis develops spectral analysis methods to detect correlations in 
		multi-dimensional time series, which are used as Early Warning Signals 
		for tipping events in stochastic dynamical systems.  
	\item Publication of three papers in respected journals:
		\begin{innerlist2}
			\item \textit{A novel scaling indicator of early warning signals} \href{https://doi.org/10.1209/0295-5075/121/10002}{(EPL, 2018)}
			\item \textit{Generalized early warning signals in multivariate and gridded data} \href{https://doi.org/10.1063/1.5093495}{(Chaos, 2019)}
			\item \textit{Power spectrum scaling as a measure of critical slowing down} \href{https://doi.org/10.1088/1748-9326/ac526f}{(ERL, 2022)}
		\end{innerlist2}
	\item Code written in \textit{MatLab} and \textit{Python}. Visualisations produced 
		using \textit{MatPlotLib} and \textit{XMGrace}.
	\item In many cases I implemented mathematically derived methods —both my own 
		and those of others— numerically by discretising and linearising. 
	\item Participation in several international conferences and workshops. 
	\item Three-month data-focussed internship at the UK Met Office (see \textit{Professional Exp.} section).
\end{innerlist}
\vspace{14pt}

\textbf{M.Res. Mathematics} — \textbf{Imperial College London} (\emph{Distinction})\\
\vspace{-2pt}
\hspace{2pt} \textbf{Dissertation:} \href{https://github.com/DrPrettyman/dissertations/blob/main/MresDiss_Prettyman.pdf}{\emph{Mesh Generation Using Solutions of an Optimal Transport Problem}}. 
\vspace{4pt}
\begin{innerlist}
	% \item Dissertation title: \href{https://github.com/DrPrettyman/dissertations/blob/main/MresDiss_Prettyman.pdf}{\emph{Adaptive Mesh Generation}}.
	\item My dissertation uses solutions to a linearisation of the Monge-Ampère equation 
		to generate an adaptive, optimally distributed mesh for numerical models.
	\item Working in C++; monitoring CPU usage and convergence time over various scenarios. 
	\item Participation in weekly soft skills workshops and monthly internal showcases.
	\item Attendance at weekly Mathematics of Planet Earth seminars 
		hosted by the Center for Doctoral Training. 
	\item Completed taught courses in Python, R, Dynamical Systems, 
		Statistic differential equations, Bayesian Analysis and Probability.
\end{innerlist}
\vspace{14pt}

\textbf{MA Mathematics} — \textbf{University of Edinburgh} (\emph{First Class Hons.}) \\
\vspace{-2pt}
\hspace{2pt} \textbf{Dissertation:} \href{https://github.com/DrPrettyman/dissertations/blob/main/MaDiss_Prettyman.pdf}{\emph{Growth and its Applications in Graph Theory}}. 
\vspace{4pt}
\begin{innerlist}
	\item My dissertation explores patterns in paths through connected digraphs. 
	\item For the dissertation project I wrote a program with a simple UI to 
		investigate the properties of these digraphs: this then became a teaching aid, 
		earning a letter of thanks from the principal. 
	\item I completed the requisite taught courses for the Pure Mathematics degree, 
		besides optional modules in computing and mathematics education.
\end{innerlist}


\end{document}

%%%%%%%%%%%%%%%%%%%%%%%%%% End CV Document %%%%%%%%%%%%%%%%%%%%%%%%%%%%%